\cotefile{156-4}
\scanpage{50}
\mspage{49}

\deleted{\uline{\unsure{Cont|Compt|Const} 5} Si $F$ et $G$ sont deux figures disjointes, et \unsure{si} $F'$, \deleted{\unsure{$G'$}\illegible{1} de} subdivision\unsure{s|{}}\unsure{,} \inserted{de $F$} alors $F'$ et $G$ sont disjointes.} \margin{\deleted{\uline{NB} $F'\cap G$ est défini \unsure{comme|com.}, \unsure{c'est} $F'$ et $G$ \uline{compatibles} qui est \unsure{en cause}.}}

\deleted{\uline{Cor} Si $F'$ \unsure{raffine|subfine} $F$, $G'$ \unsure{raffine} $G$, \deleted{\unsure{alors}} $F$ et $G$ disj. alors $F'$ et $G'$ disjointes.}

\uline{C\deleted{\unsure{ont|ompt|onst}} 5} Soient $F$ une figure, $F'$ une subdivision de $F$. \deleted{\unsure{A}}Alors pour \unsure{tt} figure $G$, $G$ est disjointe \unsure{;|:|à} $F$ ssi elle est disjoint\unsure{e} de $F'$

\uline{NB} Il suffit de le poser pour $G$ élémentaire.

Il résulte alors de \deleted{\unsure{Cont|Compt|Const} 4} C4 que $F'\cup G$ est une subdivision de $F\cup G$.

\uline{Cor} Soient $F$ et $G$ figures disjointes, alors pour $F'\ll F$, $G'\ll F$, $F'$ et $G'$ sont disjointes, et $F'\cup G'\ll F\cup G$. Si $F'$ et $G'$ sont des \unsure{subd} de $F,G$, alors $F'\cup G'$ en est une de $F\cup G$.

\uline{C 6} (Subdivisions induites) Soient $F$ une figure, $F'$ une subdivision de $F$, $G$ une sous-figure de $F$. Alors il existe une unique subdivision $G'$ de $G$, telle que $G'\subset F'$ (subdivision induite \inserted{et $G'$ est formée des $X\in F'$ telles que $X\ll G$})

\margin{NB \inserted{Double} Transitivité de l'induction. \uline{Notation} $F'|G$}

\uline{C 7} (\unsure{Recollement|Récollement} des subdivisions) Soit $F$ une figure, \deleted{\unsure{chaque} \illegible{3} \unsure{tt}} \inserted{$F_i$ des sous-figures} \deleted{$X\in$} telles que $F=\bigcup F_i$. Soit, pour \unsure{tt} $i$, $F'_i$ une subdivision de $F_i$, et supposons
